\begin{myChapterIntro}{本章涵盖}
\begin{itemize}
\item 用递归为编程问题设计解决方案
\item 动态回溯，以在运行时尝试不同的求解路径
\end{itemize}
\end{myChapterIntro}

递归（recursion）是一种软件设计技术，用于为某些编程问题设计解决方案。它涉及一个调用自身的函数，几乎所有现代编程语言都支持这一重要技术。如果我们对某些类别的编程问题恰当地使用递归，就能设计出简单、优雅，甚至有人会称之为神奇的解决方案。借助递归，我们可以设计出用其他方式难以求解的方案。

\myGraphic{0.893}{images/CH15_00_UN01_Mak.png}{}

递归要求我们换一种思路来思考如何为某个问题设计编程解决方案。遗憾的是，对某些程序员来说，递归太过神秘、令人望而却步，以至于不敢使用。本章将拨开这层迷雾，并演示如何把递归与动态回溯结合起来，从而获得更强大的设计工具。不过，我们也会看到误用递归如何导致出人意料的性能灾难。随着经验的积累，我们能学会何时适合用递归来设计。

\section{递归与 for 循环的比较}

通过把递归与我们熟悉的 \texttt{for} 循环作比较，可以消除笼罩在递归之上的部分神秘感。递归背后的核心思想是：把一个编程问题分解为一个规模更小但性质相似的问题，再用同样的方法求解。解决较小的问题，就能进而解决较大的问题。我们不断把问题分解为更小但相似的问题，直到问题小到其解显而易见、唾手可得。这个解会成为次小问题之解的一部分，而后者又成为下一个更大问题之解的一部分。当递归展开、回到最大的问题时，那些较小问题的解就成了原问题之解的一部分。

我们编写一个递归函数来求解某个编程问题，该函数通过调用自身来求解一个规模更小但相似的问题。函数调用自身时就构成一次递归调用，通常称为\emph{嵌套}调用。递归调用不断重复、不断向内嵌套，以求解越来越小的问题，直到问题小到函数无需再次递归调用就能立即返回结果，从而避免无限递归。这个最小的问题称为递归的\emph{基例}（base case）。基例返回的结果使倒数第二次递归调用得以带着结果返回，随后所有递归调用的结果沿着调用链一个接一个地返回到最初的调用。最初那次函数调用的返回就解决了这个编程问题。正是这些递归调用以及调用展开时返回的结果，让此类解决方案变得简单而优雅。

理解递归的一种方法，是把它与下面这种基本的计数 \texttt{for} 循环作比较：

\begin{codebox}
for (int i = INITIAL_VALUE; i < LIMIT_VALUE; i++) ...
\end{codebox}

表 15.1 展示了这一比较。

\begin{center}
{\bfseries 表 15.1 把 \texttt{for} 循环与递归作比较}\par\vspace{3bp}
\begin{tabularx}{\textwidth}{XXX}
\toprule
 & \texttt{for} 循环 & 递归 \\
初始条件 & \texttt{i =\allowbreak{} INITIAL\_\allowbreak{}VALUE} & 原始问题 \\
重复更新 & \texttt{i++} & 不断把问题分解为更小但相似的问题。 \\
终止条件 & \texttt{i >=\allowbreak{} LIMIT\_\allowbreak{}VALUE} & 问题小到其解显而易见、唾手可得（一个基例）。 \\
\bottomrule
\end{tabularx}
\end{center}

\texttt{for} 循环最终能够结束，这一点至关重要。因此，循环的每一次迭代都必须更新控制变量，使其取值逐步逼近终止条件。同样，要让递归最终结束，问题必须变得越来越小并趋近某个基例；在基例处，问题小到其解显而易见、唾手可得，不再有任何递归调用。

15.2 节和 15.3 节中的两个示例程序演示了递归。这两个程序所解决的问题都不应该用递归来做——用 \texttt{for} 或 \texttt{while} 循环这类常规迭代来解效率要高得多。但我们会用它们来说明递归是如何工作的。

\section{用递归求最大值}

我们的第一个示例程序用递归来求一个整数 vector 中的最大值。要设计这个递归解，我们必须确定表 15.1 最后一列的各个条目：

\begin{itemize}
\item \emph{初始条件}——原始的取值 vector。
\item \emph{更新}——一个规模更小但相似的问题：在更短的 vector 中求最大值。
\item \emph{终止条件}——基例：vector 只包含单个元素（或根本不包含任何元素）。
\end{itemize}

图 15.1 展示了递归的工作过程。我们首先把原始取值 vector 传给递归函数 \texttt{largest()}。该函数把 vector 第一个元素的值记录到变量 \texttt{first\_\allowbreak{}value} 中，然后移除第一个元素，得到一个更短的 vector，即原 vector 的其余部分。函数把 \texttt{first\_\allowbreak{}value} 与其余部分中的最大值作比较，两者中较大者就是原问题的解。

可是函数 \texttt{largest()} 怎么知道 vector 其余部分中的最大值呢？那正是以更短 vector 递归调用该函数所得的结果。每次递归调用接收到的 vector 实参都比原先少一个元素，因此这些调用逐步逼近基例。每次递归调用返回其问题的解。最后一次调用收到的 vector 实参只包含一个元素，这就是该调用立即、显而易见的返回值。最后一次调用不再发起递归调用，从而终止了递归调用链。这个立即返回的值使基例的调用者得以完成比较并返回一个值。随着递归展开，每次递归调用的返回值都使调用者得以完成比较并返回结果，直到最初的调用返回整个 vector 的最大值。

\myGraphic{0.980}{images/CH15_F01_Mak.png}{图 15.1 递归函数 \texttt{largest()} 如何求 vector 中的最大值。加框的取值是各次调用中不断变短的 vector 实参。每次递归调用的返回值都使调用者得以返回其值。}

递归函数 \texttt{largest()} 的工作方式正如图 15.1 所示。

\begin{codeboxt}{代码清单 15.1（程序 15.1 Largest）：\texttt{largest.\allowbreak{}cpp}（仅用于演示递归）}
#include <iostream>
#include <iomanip>
#include <vector>
#include <stdlib.h>
#include <time.h>

using namespace std;

const int SIZE = 12;

int largest(vector<int> v);

int main()
{
    vector<int> data;

    srand(time(0));
    for (int i = 0; i < SIZE; i++) data.push_back(rand()%100);

    cout << "Largest of";
    for (int d : data) cout << setw(3) << d;

    cout << endl;
    cout << "is" << setw(3) << largest(data);               ①
    return 0;
}

int largest(vector<int> v)                                  ②
{
    if (v.size() == 1) return v[0];                         ③

    int first_value = v[0];                                 ④
    v.erase(v.begin());                                     ⑤

    int largest_of_rest = largest(v);                       ⑥

    return first_value > largest_of_rest ? first_value      ⑦
                                         : largest_of_rest; ⑦

}
\end{codeboxt}
\noindent{\fontsize{9bp}{12bp}\selectfont ① 传入原始 vector 的初始函数调用}\par
\noindent{\fontsize{9bp}{12bp}\selectfont ② 递归函数}\par
\noindent{\fontsize{9bp}{12bp}\selectfont ③ 大小为 1 的 vector 属于基例，无需递归调用}\par
\noindent{\fontsize{9bp}{12bp}\selectfont ④ 记住第一个元素的值。}\par
\noindent{\fontsize{9bp}{12bp}\selectfont ⑤ 移除第一个元素以缩短 vector。}\par
\noindent{\fontsize{9bp}{12bp}\selectfont ⑥ 递归调用，求 vector 其余部分中的最大值}\par
\noindent{\fontsize{9bp}{12bp}\selectfont ⑦ 把第一个元素的值与 vector 其余部分中的最大值作比较。}\par

以下是用随机生成的取值 vector 运行时得到的示例输出：

\begin{codebox}
Largest of 84 87 60 50 81 15 68 82 92  1 59 31
is 92
\end{codebox}

\myGraphic{0.811}{images/CH15_F01_UN02_Mak.png}{}

\section{用递归反转 vector}

再来看一个示例程序，它展示了递归的工作方式。我们将编写一个递归函数 \texttt{reverse()} 来反转 vector 的内容。与上一个例子一样，这并不是解决该问题的最佳方式，但它是展示递归工作原理的又一个好例子。

我们首先把原始 vector（例如包含整数 10、20、30、40、50）传给递归函数 \texttt{reverse()}。该函数移除 vector 的第一个元素，递归地反转其余部分（更短，如下面方括号内所示），然后把被移除的第一个元素追加到 vector 末尾：

\begin{codebox}
10 [ 20 30 40 50 ]
   [ 50 40 30 20 ] 10
\end{codebox}

我们把这个更短的 vector 传给 \texttt{reverse()} 的一次递归调用来反转它。每次递归调用得到的 vector 都更短。当 vector 最终到达只含一个元素的基例时，递归结束。最终结果是一个包含 50、40、30、20、10 的 vector。图 15.2 展示了如何递归地反转 vector 的内容。

\myGraphic{0.918}{images/CH15_F02_Mak.png}{图 15.2 由递归函数 \texttt{reverse()} 反转的 vector 内容。每次调用都移除第一个值，递归地反转 vector 其余部分，然后把被移除的第一个值追加到末尾。}

递归函数 \texttt{reverse()} 的工作方式正如图 15.2 所示。

\begin{codeboxt}{代码清单 15.2（程序 15.2 Reverse）：\texttt{reverse.\allowbreak{}cpp}（仅用于演示递归）}
#include <iostream>
#include <iomanip>
#include <vector>

using namespace std;

const int SIZE = 5;

void reverse(vector<int>& v);

int main()
{
    vector<int> data;

    for (int i = 1; i <= SIZE; i++) data.push_back(10*i);

    cout << "Reverse of";
    for (int d : data) cout << setw(3) << d;
    cout << endl;

    reverse(data);                        ①

    cout << "        is";
    for (int d : data) cout << setw(3) << d;
    cout << endl;

    return 0;
}

void reverse(vector<int>& v)               ②
{
    if (v.size() == 1) return;             ③

    int first_value = v[0];                ④
    v.erase(v.begin());                    ⑤

    reverse(v);                            ⑥
    v.push_back(first_value);              ⑦
}
\end{codeboxt}
\noindent{\fontsize{9bp}{12bp}\selectfont ① 初始函数调用}\par
\noindent{\fontsize{9bp}{12bp}\selectfont ② 递归函数}\par
\noindent{\fontsize{9bp}{12bp}\selectfont ③ 大小为 1 的 vector 属于基例，无需递归调用}\par
\noindent{\fontsize{9bp}{12bp}\selectfont ④ 记住第一个元素的值。}\par
\noindent{\fontsize{9bp}{12bp}\selectfont ⑤ 移除第一个元素以缩短 vector。}\par
\noindent{\fontsize{9bp}{12bp}\selectfont ⑥ 递归调用，反转 vector 其余部分的内容}\par
\noindent{\fontsize{9bp}{12bp}\selectfont ⑦ 把第一个元素追加到反转后的 vector 末尾。}\par

示例输出为

\begin{codebox}
Reverse of 10 20 30 40 50
        is 50 40 30 20 10
\end{codebox}

\myGraphic{0.728}{images/CH15_F02_UN03_Mak.png}{}

\section{用递归求解汉诺塔难题}

汉诺塔（Towers of Hanoi）难题为展示递归的恰当用法提供了绝佳机会。这个问题有一个非常优雅的递归解，但不用递归则很难求解。我们需要运用递归思维。

图 15.3 展示了四圆盘难题的起始位置与结束位置。每个圆盘中心都有一个孔，它们按大小顺序套在三根柱子之一上，最小的圆盘在最上面。目标是把所有圆盘从源柱移到目标柱，并在那里重新按大小顺序叠放。我们给柱子标上 L、M、R（分别代表左、中、右）。

\myGraphic{0.296}{images/CH15_F03_Mak.png}{图 15.3 汉诺塔难题中四个圆盘的起始位置与结束位置。我们把这些圆盘从源柱 L 移到了目标柱 R。}

移动圆盘的规则如下：

\begin{enumerate}
\item 每次只能移动一个圆盘。
\item 绝不能把较大的圆盘放到较小的圆盘上面。
\item 用第三根柱子作为移动圆盘时的临时落脚点。
\end{enumerate}

图 15.4 从概念上展示了如果不存在第 1 条规则我们会如何求解该难题。我们设法把上面三个圆盘移到临时柱 M。然后把最大的圆盘从柱 L 移到柱 R。最后把三个圆盘从柱 M 移到柱 R，就完成了。

\myGraphic{0.312}{images/CH15_F04_Mak.png}{图 15.4 以柱 M 作为临时柱的四圆盘概念性解法。}

但是，如何在不违反“每次只能移动一个圆盘”这条规则的前提下移动三个圆盘呢？移动三个圆盘是比移动四个圆盘规模更小的问题。这个问题提示了递归解法。

为了确信递归解法可行，我们先从基例入手。图 15.5 表明，把单个圆盘从源柱 L 移到目标柱 R 是一次直接操作——不需要递归。因此，基例确实会终止进一步的递归。

\myGraphic{0.312}{images/CH15_F05_Mak.png}{图 15.5 把单个圆盘从柱 L 移到柱 R 的基例。不需要递归。}

图 15.6 展示了把两个圆盘从源柱 L 移到目标柱 R 的移动序列。这需要把柱 M 用作临时柱。共有三次单盘移动：首先从柱 L 到柱 M（L 是源，M 是这次移动的目标），然后从柱 L 到柱 R（L 是源，R 是这次移动的目标），最后从柱 M 到柱 R（M 是源，R 是这次移动的目标）。注意柱子可以变换角色——在第 1 步中 M 是目标，而在第 3 步中 M 是源。我们用圆盘数更少（只有一个）的递归来求解这些单盘移动。

\myGraphic{0.312}{images/CH15_F06_Mak.png}{图 15.6 以柱 M 作为临时柱，把两个圆盘从柱 L 移到柱 R 的移动序列}

我们可以求解三圆盘的难题。为了更清楚地演示我们在使用递归，图 15.7 求解了该难题，把它们从柱 L 移到柱 R。

\myGraphic{0.679}{images/CH15_F07_Mak.png}{图 15.7 把三个圆盘从柱 L 移到柱 R 的移动序列。在这些步骤中，各柱子会轮换充当源柱、目标柱和临时柱的角色。}

第 1、2、3 步把上面两个圆盘从柱 L 移到柱 M。我们知道可以递归地求解移动两个圆盘的问题。第 4 步是把一个圆盘从柱 L 移到柱 R 的基例。第 5、6、7 步递归地求解把两个圆盘从柱 M 移到柱 R 的问题。这就解出了三圆盘的难题，其中涉及递归地求解一个和两个圆盘的移动：

\begin{itemize}
\item 递归地求解把两个圆盘从柱 L 移到柱 M。
\item 递归地求解把一个圆盘从柱 L 移到柱 R（基例）。
\item 递归地求解把两个圆盘从柱 M 移到柱 R。
\end{itemize}

现在我们知道可以求解四圆盘的难题了——递归地把三个圆盘从柱 L 移到柱 M，递归地把一个圆盘从柱 L 移到柱 R（基例），再递归地把三个圆盘从柱 M 移到柱 R。事实上，要解出 \emph{N} 个圆盘的难题，我们递归地利用 \emph{N}-1 个圆盘和一个圆盘的解法。

\begin{codeboxt}{代码清单 15.3（程序 15.3 Towers）：\texttt{towers.\allowbreak{}cpp}}
#include <iostream>
#include <iomanip>
#include <vector>

using namespace std;

const int N = 3;

void solve(const int  n,
           const char source, 
           const char temporary,
           const char destination )

int main()
{
    solve(N, 'L', 'M', 'R');                                         ①
    return 0;
}

void solve(const int  n,                                             ②
           const char source, 
           const char temporary,
           const char destination )
{
    if (n == 1)                                                      ③
    {                                                                ③
        cout << «Move « << source << « ==> « << destination << endl; ③
    }                                                                ③
    else
    {
        solve(n - 1, source,    destination, temporary);             ④
        solve(1,     source,    temporary,   destination);           ④
        solve(n - 1, temporary, source,      destination);           ④
    }
}
\end{codeboxt}
\noindent{\fontsize{9bp}{12bp}\selectfont ① 传入 N 个圆盘及柱名与角色的初始调用}\par
\noindent{\fontsize{9bp}{12bp}\selectfont ② 递归函数}\par
\noindent{\fontsize{9bp}{12bp}\selectfont ③ 基例}\par
\noindent{\fontsize{9bp}{12bp}\selectfont ④ 用更少圆盘进行的递归调用}\par

以下是把 \texttt{N = 3} 个圆盘求解时的输出：

\begin{codebox}
Move L ==> R
Move L ==> M
Move R ==> M
Move L ==> R
Move M ==> L
Move M ==> R
Move L ==> R
\end{codebox}
\noindent{\fontsize{9bp}{12bp}\selectfont ① 传入 N 个圆盘及柱名与角色的初始调用}\par
\noindent{\fontsize{9bp}{12bp}\selectfont ② 递归函数}\par
\noindent{\fontsize{9bp}{12bp}\selectfont ③ 基例}\par
\noindent{\fontsize{9bp}{12bp}\selectfont ④ 用更少圆盘进行的递归调用}\par

\myGraphic{0.776}{images/CH15_F07_UN04_Mak.png}{}

\section{二叉搜索树的递归算法}

对于递归定义的数据结构来说，递归算法可能再合适不过了。二叉树就是这样一种递归数据结构，我们将看到如何设计递归函数来操作树中的节点。

\emph{二叉树}是一种树形结构，其中每个非空节点有零个、一个或两个子节点。因此，节点有两条链接：左子链接和右子链接。每条链接要么是 \texttt{nullptr}，要么指向某个子节点，而该子节点是某棵子树的根（图 15.8）。每棵子树本身也是一棵二叉树，因此二叉树是递归定义的。根代表整棵树。我们可以把 \texttt{nullptr} 链接看作指向一个空节点的指针，该空节点代表一棵空树。越往下走子树越小，直到到达空树这一基例。

\myGraphic{0.705}{images/CH15_F08_Mak.png}{图 15.8 一棵二叉树，其中每个非空节点有零个、一个或两个子节点，这些子节点是各子树的根。每棵子树本身也是一棵二叉树，因此二叉树是递归定义的。这棵二叉树同时还是二叉搜索树（binary search tree，BST）：在每个非空父节点处，左子树中的值都小于或等于父节点的值，而右子树中的值都更大。每棵子树本身也是一棵 BST，因此 BST 是递归定义的。}

如果一棵二叉树还满足以下条件，它就是\emph{二叉搜索树}（BST）：

\begin{itemize}
\item 每个非空节点都有一个值。
\item 任何非空父节点的左子节点（如果存在）是一棵子树，其所有节点的值都小于或等于父节点的值。
\item 任何非空父节点的右子节点（如果存在）是一棵子树，其所有节点的值都大于父节点的值。
\end{itemize}

BST 是递归定义的，因为每个节点都是一棵子树的根，而这棵子树本身又是一棵 BST。

\begin{codeboxt}{代码清单 15.4（程序 15.4-BST）：\texttt{Node.h}}
#include "BST.h"

class Node
{
public:
    Node(const int v) : value(v), left(nullptr), right(nullptr) {}
    friend class BST;

public:
    int value;
    Node *left;                                         ①
    Node *right;                                        ①
};
\end{codeboxt}
\noindent{\fontsize{9bp}{12bp}\selectfont ① 指向左、右子节点的指针，每个子节点都是一棵 BST 的根}\par

\subsection{15.5.1 用递归向 BST 中插入节点}

我们将为 BST 定义三个递归操作：插入节点、打印所有节点、删除所有节点。两个公有成员函数 \texttt{insert()} 和 \texttt{print()} 从树的根开始，调用相应的私有成员函数。

\begin{codeboxt}{代码清单 15.5（程序 15.4-BST）：\texttt{BST.h}}
#include "Node.h"

class BST
{
public:
    BST() : root(nullptr) {}

    virtual ~BST() { remove(root); }                      ①

    void insert(const int value) { insert(value, root); } ②
    void print() const           { print(root); }         ③

private:
    Node *root;

    void insert(const int value, Node*& link);            ④
    void print (const Node * const link) const;           ④
    void remove(const Node * const link) const;           ④
};
\end{codeboxt}
\noindent{\fontsize{9bp}{12bp}\selectfont ① 从根开始删除。}\par
\noindent{\fontsize{9bp}{12bp}\selectfont ② 从根开始插入。}\par
\noindent{\fontsize{9bp}{12bp}\selectfont ③ 从根开始打印。}\par
\noindent{\fontsize{9bp}{12bp}\selectfont ④ 递归成员函数}\par

我们向私有成员函数 \texttt{insert()}、\texttt{print()} 和 \texttt{remove()} 各传入一个指向节点的指针。该指针可以是 \texttt{nullptr}。我们以引用方式向成员函数 \texttt{insert()} 传入该指针，以便该函数能够设置指针。我们还传入要插入树中的值。私有成员函数 \texttt{insert()}、\texttt{print()} 和 \texttt{remove()} 都采用递归方式工作。

\begin{codeboxt}{代码清单 15.6（程序 15.4-BST）：\texttt{BST.cpp}（第 1 部分，共 3 部分）}
#include <iostream>
#include <iomanip>
#include "BST.h"

using namespace std;

void BST::insert(const int value, Node*& link)
{
    if      (link == nullptr)      link = new Node(value);     ①
    else if (value <= link->value) insert(value, link->left);  ②
    else                           insert(value, link->right); ③
}
...
\end{codeboxt}
\noindent{\fontsize{9bp}{12bp}\selectfont ① 创建新节点并在此处插入。}\par
\noindent{\fontsize{9bp}{12bp}\selectfont ② 否则，插入到左子树中的某个位置。}\par
\noindent{\fontsize{9bp}{12bp}\selectfont ③ 或者，插入到右子树中的某个位置。}\par

我们向私有成员函数 \texttt{insert()} 传入要插入树中的新值，并以引用方式传入一个指向节点的指针。该函数首先检查指针是否为 \texttt{nullptr}，即空子树这一基例。如果是，函数就用该值创建一个新节点，并把指针设为这个新节点，从而替换空子树并插入新值。

否则，如果我们传入的指针指向一个已存在的节点，私有 \texttt{insert()} 成员函数就递归地调用自身，根据新值与节点值的比较结果，把新值插入该节点的左子树或右子树中。每次调用都作用于一棵更小的子树，更接近空子树这一基例。在基例处，一个包含新值的节点会替换该空子树，从而把新值插入到它在树中的正确位置（图 15.9）。由于公有 \texttt{insert()} 成员函数传入的是指向树根的指针（代码清单 15.5），当树最初为空时，私有 \texttt{insert()} 成员函数会把根指针设为一个包含第一个值的新节点。

\myGraphic{0.871}{images/CH15_F09_Mak.png}{图 15.9 这些值按 42、30、5、35、63、75、25、51、44、19、55 的顺序插入这棵 BST。为了把新值 60 插入树中，私有函数 \texttt{insert()} 被递归地调用，直到到达空子树这一基例。新节点替换该空子树，值就这样被插入到树中的正确位置。}

\subsection{15.5.2 用递归打印 BST}

私有成员函数 \texttt{print()} 按有序顺序递归地打印 BST 各节点的值（图 15.10）。

\myGraphic{0.748}{images/CH15_F10_Mak.png}{图 15.10 按有序顺序递归地打印 BST 的值。在每个父节点处，先打印其左子树中的值。接着打印父节点的值，然后打印父节点右子树中的值。每个节点旁的小数字表示节点的打印顺序。}

我们向私有成员函数 \texttt{print()} 传入一个指向要打印的节点的指针（代码清单 15.7）。公有 \texttt{print()} 函数通过把根指针传给私有函数来启动打印（代码清单 15.5）：

\begin{enumerate}
\item 函数首先递归地调用自身，传入父节点的左子链接，以打印小于或等于父节点值的那些值。
\item 接着，函数打印父节点的值。
\item 然后，函数再次递归地调用自身，这次传入父节点的右子链接，以打印大于父节点值的那些值。
\end{enumerate}

\begin{codeboxt}{代码清单 15.7（程序 15.4-BST）：\texttt{BST.cpp}（第 2 部分，共 3 部分）}
...

void BST::print(const Node * const link) const
{
    if (link != nullptr)
    {
        print(link->left);                 ①
        cout << setw(3) << link->value;    ②
        print(link->right);                ③
    }
}

...
\end{codeboxt}
\noindent{\fontsize{9bp}{12bp}\selectfont ① 递归调用，打印小于或等于该节点值的那些值。}\par
\noindent{\fontsize{9bp}{12bp}\selectfont ② 打印该节点的值。}\par
\noindent{\fontsize{9bp}{12bp}\selectfont ③ 递归调用，打印大于该节点值的那些值。}\par

\subsection{15.5.3 用递归删除 BST 的所有节点}

递归地删除节点与打印节点类似，只不过由析构函数调用的私有成员函数 \texttt{remove()}（代码清单 15.5）会先递归地删除一个节点的两棵子树，然后再删除该节点本身。

\begin{codeboxt}{代码清单 15.8（程序 15.4-BST）：\texttt{BST.cpp}（第 3 部分，共 3 部分）}
...

void BST::remove(const Node * const link) const
{
    if (link != nullptr)
    {
        remove(link->left);                 ①
        remove(link->right);                ①

        delete link;                        ②
    }
}
\end{codeboxt}
\noindent{\fontsize{9bp}{12bp}\selectfont ① 递归地删除两棵子树。}\par
\noindent{\fontsize{9bp}{12bp}\selectfont ② 删除该节点本身}\par

测试程序生成随机整数值，把它们插入 BST，并按有序顺序打印树中的值。

\begin{codeboxt}{代码清单 15.9（程序 15.4-BST）：\texttt{tester.\allowbreak{}cpp}}
#include <iostream>
#include <iomanip>
#include <stdlib.h>
#include <time.h>
#include "BST.h"

using namespace std;

int main()
{
    const int TREE_SIZE = 20;

    BST tree;
    srand(time(0));

    cout << "Inserting:";
    for (int i = 0; i < TREE_SIZE; i++)
    {
        int value = rand()%100;
        tree.insert(value);
        cout << setw(3) << value;
    }

    cout << endl << " Printing:";
    tree.print();
}
\end{codeboxt}

以下是示例程序输出：

\begin{codebox}
Inserting: 75 56 60 38 22  9 98 58 10 98 30 55 27 58 95 42 71 80 29 52
 Printing:  9 10 22 27 29 30 38 42 52 55 56 58 58 60 71 75 80 95 98 98
\end{codebox}

\begin{mySidebar}{树的遍历}

递归函数 \texttt{print()} 执行的是\emph{中序}（inorder）树遍历——函数在访问节点的左子树和右子树\emph{之间}访问该节点（以打印节点的值）。递归函数 \texttt{remove()} 执行的是\emph{后序}（postorder）树遍历——函数在访问节点的左子树和右子树\emph{之后}访问该节点（以删除该节点）。
\end{mySidebar}

\myGraphic{0.790}{images/CH15_F10_UN05_Mak.png}{}

\section{用递归对数组进行快速排序}

快速排序（quicksort）是一种快速对取值数组进行排序的算法。它是递归的又一个绝佳应用示例。

快速排序依赖一个把数组划分成两个子数组的过程。例如，要对下面这个数组进行快速排序

\myGraphic{0.499}{images/CH15_F10_UN06_Mak.png}{}

我们首先选取数组中位于中央或接近中央的一个元素作为\emph{基准元素}（pivot element）：16（如上文圆圈所示）。按照下文说明的过程，我们重新排列各元素，使所有小于或等于基准值的值都位于基准的左侧，所有大于基准值的值都位于基准的右侧：

\myGraphic{0.510}{images/CH15_F10_UN07_Mak.png}{}

现在，我们在基准两侧得到两个独立的子数组（如上文方框所示）。这两个子数组都未必有序，我们将对它们分别递归地进行快速排序。

快速排序算法为：

\begin{enumerate}
\item 选取一个基准元素。（在本例中，我们将选取数组中位于中央或接近中央的元素。）
\item 把待排序数组划分成位于基准两侧的两个子数组。 a. 基例 1：如果子数组的大小为 0 或 1，则不做任何事，因为它已经有序。 b. 基例 2：如果子数组的大小为 2，则在必要时交换这两个元素。 c. 否则，对这两个子数组递归地进行快速排序。
\end{enumerate}

关于该算法的一些关键事实如下：

\begin{itemize}
\item 在选取基准元素并完成划分之后，基准元素在最终的有序数组中已经处于正确的位置。
\item 满足任一基例的划分中的元素（在基例 2 中必要时完成交换之后）在最终的有序数组中已经处于正确的位置。
\end{itemize}

\subsection{15.6.1 快速排序实战}

图 15.11 展示了快速排序算法对数组排序时的实际过程。第 A 行是原始数组，第 L 行是有序数组。随着排序的推进，带圆圈的值是每一步的基准值，加粗的值是那些在最终有序数组中已经处于正确位置的值。需要递归地快速排序的子数组用方框标出，每一步中正在被快速排序的子数组用阴影表示。

\myGraphic{0.707}{images/CH15_F11_Mak.png}{图 15.11 快速排序算法对数组排序时的实际过程。带圆圈的值是被选中的基准元素，它们在划分后于两侧形成子数组。矩形框住的子数组需要递归地进行快速排序。在每一行中，带阴影的子数组是当前正在被排序的那个。加粗的值在最终有序数组中已经处于正确位置。注意，被选中的基准元素在划分后总是已经处于正确位置。}

第 B 行表明我们开始对整个数组排序。我们用基准元素 7 把数组划分成两个子数组（第 C 行）。我们递归地对由 4 和 3 这两个值组成的第一个子数组排序，交换了它们。这是基例 2。

接下来（第 D 行），我们递归地对另一个子数组排序，首先选取 71 作为基准元素。这又产生了两个子数组（第 E 行）。我们递归地对左子数组排序，在其中选取 65 作为基准元素（第 E 行）。由于划分后基准值落到了子数组的右端，65 右侧得到的子数组为空（基例 1）。

第 G、H、J、K 行展示的划分都只包含一个元素（基例 1）。它们的值已经处于正确位置。

\myGraphic{0.664}{images/CH15_F11_UN08_Mak.png}{}

\subsection{15.6.2 把数组划分成子数组}

快速排序算法中的关键过程是把数组划分成两个子数组。图 15.12 展示了如何以基准值 52（带圆圈）来完成划分。变量 \texttt{i} 和 \texttt{j} 是数组下标。划分完成后，52 左侧的所有值都更小，52 右侧的所有值都更大。这两个（加框的）子数组现在可以递归地进行快速排序了。

\myGraphic{0.793}{images/CH15_F12_Mak.png}{图 15.12 选取 52（带圆圈）作为基准值，然后把数组划分成两个子数组。变量 \texttt{i} 和 \texttt{j} 是数组下标。划分之后，52 左侧的所有值都小于或等于它，52 右侧的所有值都大于它。基准两侧这两个加框的子数组现在可以递归地进行快速排序了。}

划分数组的步骤如下：

\begin{enumerate}
\item 选取数组中的一个元素作为基准元素。在我们的示例中，将选取位于中央或接近中央的元素（第 A 行）。
\item 把基准元素与数组最右边的元素交换，将它暂时挪开（第 B 行）。把变量 \texttt{i} 设为最左边元素的下标，把变量 \texttt{j} 设为紧挨在被挪开的基准元素左侧那个元素的下标。
\item 当第 \texttt{i} 个元素的值小于或等于基准值时，把 \texttt{i} 向右移动（即递增 \texttt{i}）。（如果该条件一开始就不成立，它可能完全不动。）换句话说，\texttt{i} 会跳过那些已经正确位于基准值左侧的值。
\item 当第 \texttt{j} 个元素的值大于基准值时，把 \texttt{j} 向左移动（即递减 \texttt{j}）（第 C 行）。（如果该条件一开始就不成立，它可能完全不动。）换句话说，\texttt{j} 会跳过那些已经正确位于基准值右侧的值。
\item 在 \texttt{i} 和 \texttt{j} 都停止移动之后，如果 \texttt{i < j}，则交换第 \texttt{i} 个和第 \texttt{j} 个元素的值（第 D 行）。
\item 重复第 3、4、5 步（第 E 行到第 I 行），直到 \texttt{j} 越过 \texttt{i}（第 I 行）。把第 \texttt{i} 个元素的值与我们此前挪到右边的基准元素交换（第 J 行）。此时数组已划分成位于基准元素两侧的两个子数组。
\end{enumerate}

下面是一次包含所有细节的完整快速排序。每个基准值都用圆括号标出，当前正在被递归快速排序的子数组用方括号括起：

\begin{codebox}
  86   95   51   95   28   33    9   15   84   86   67 
-------------------------------------------------------
 [86   95   51   95   28  (33)   9   15   84   86   67]  pivot (33)
 [86   95   51   95   28   67    9   15   84   86   33]  67 <=> 33
   i                                            j      
 [86   95   51   95   28   67    9   15   84   86   33]  moved j
   i                                  j                
 [15   95   51   95   28   67    9   86   84   86   33]  15 <=> 86
   i                                  j                
 [15   95   51   95   28   67    9   86   84   86   33]  moved i
        i                             j                
 [15   95   51   95   28   67    9   86   84   86   33]  moved j
        i                        j                     
 [15    9   51   95   28   67   95   86   84   86   33]  9 <=> 95
        i                        j                     
 [15    9   51   95   28   67   95   86   84   86   33]  moved i
             i                   j                     
 [15    9   51   95   28   67   95   86   84   86   33]  moved j
             i         j                               
 [15    9   28   95   51   67   95   86   84   86   33]  28 <=> 51
             i         j                               
 [15    9   28   95   51   67   95   86   84   86   33]  moved i
                  i    j                               
 [15    9   28   95   51   67   95   86   84   86   33]  moved j
             j    i                                    
 [15    9   28  (33)  51   67   95   86   84   86   95]  33 <=> 95
-------------------------------------------------------
 [15  ( 9)  28]  33   51   67   95   86   84   86   95   pivot (9)
 [15   28    9]  33   51   67   95   86   84   86   95   28 <=> 9
   i    j                                              
 [15   28    9]  33   51   67   95   86   84   86   95   moved j
   i                                                   
 [ 9)  28   15]  33   51   67   95   86   84   86   95    9 <=> 15
-------------------------------------------------------
   9  [28   15]  33   51   67   95   86   84   86   95 
   9  [15   28]  33   51   67   95   86   84   86   95   15 <=> 28
-------------------------------------------------------
   9   15   28   33  [51   67   95  (86)  84   86   95]  pivot (86)
   9   15   28   33  [51   67   95   95   84   86   86]  95 <=> 86
                       i                        j      
   9   15   28   33  [51   67   95   95   84   86   86]  moved i
                                 i                   j 
   9   15   28   33  [51   67   86   95   84   95   86]  86 <=> 95
                                 i              j      
   9   15   28   33  [51   67   86   95   84   95   86]  moved i
                                      i         j      
   9   15   28   33  [51   67   86   95   84   95   86]  moved j
                                      i    j           
   9   15   28   33  [51   67   86   84   95   95   86]  84 <=> 95
                                      i    j           
   9   15   28   33  [51   67   86   84   95   95   86]  moved i
                                          ji           
   9   15   28   33  [51   67   86   84   95   95   86]  moved j
                                      j    i           
   9   15   28   33  [51   67   86   84  (86)  95   95]  86 <=> 95
-------------------------------------------------------
   9   15   28   33  [51  (67)  86   84]  86   95   95   pivot (67)
   9   15   28   33  [51   84   86   67]  86   95   95   84 <=> 67
                       i         j                     
   9   15   28   33  [51   84   86   67]  86   95   95   moved i
                            i         j                
   9   15   28   33  [51   84   86   67]  86   95   95   moved j
                       j    i                          
   9   15   28   33  [51  (67)  86   84]  86   95   95   67 <=> 84
-------------------------------------------------------
   9   15   28   33  [51]  67   86   84   86   95   95 
-------------------------------------------------------
   9   15   28   33   51   67  [86   84]  86   95   95 
   9   15   28   33   51   67  [84   86]  86   95   95   84 <=> 86
-------------------------------------------------------
   9   15   28   33   51   67   84   86   86  [95   95]
\end{codebox}

\subsection{15.6.3 快速排序的实现}

类 \texttt{Quicksort} 实现了快速排序算法，如下面的代码清单所示。重要的成员函数有递归的 \texttt{sort()} 函数，以及 \texttt{partition(\allowbreak{})\allowbreak{}} 和 \texttt{swap\_\allowbreak{}values(\allowbreak{})\allowbreak{}} 函数。

\begin{codeboxt}{代码清单 15.10（程序 15.5 Quicksort）：\texttt{Quicksort.\allowbreak{}h}}
#include <iostream>
#include <iomanip>

using namespace std;

class Quicksort
{
public:
    Quicksort(int * const d, const int s) : data(d), size(s) {};

    void sort() { sort(0, size-1); }

    bool verify_sorted();

private:
    int *data;                                                 ①
    int size;

    void sort(const int left_index, const int right_index);    ②
    int partition(const int left_index, const int right_index);
    void swap_values_at(const int index1, const int index2);

    friend ostream& operator <<(ostream& ostr, const Quicksort& q);
};
\end{codeboxt}
\noindent{\fontsize{9bp}{12bp}\selectfont ① 待排序的数据}\par
\noindent{\fontsize{9bp}{12bp}\selectfont ② 递归排序函数}\par

成员函数 \texttt{sort()} 实现基例，以及划分后对两个子数组的递归调用。待排序子数组最左边和最右边元素的下标分别是 \texttt{left\_\allowbreak{}index} 和 \texttt{right\_\allowbreak{}index}。

\begin{codeboxt}{代码清单 15.11（程序 15.5 Quicksort）：\texttt{Quicksort.\allowbreak{}cpp}（第 1 部分，共 3 部分）}
#include <iostream>
#include <iomanip>
#include "Quicksort.h"

using namespace std;

void Quicksort::sort(const int left_index, const int right_index)
{
    int partition_size = right_index - left_index + 1;

    if (partition_size < 2) return;                              ①

    if (partition_size == 2)                                     ②
    {
        if (data[left_index] > data[right_index])
        {
            swap_values_at(left_index, right_index);
        }
    }
    else
    {
        int pivot_index = partition(left_index, right_index);    ③

        sort(left_index, pivot_index - 1);                       ④
        sort(pivot_index + 1, right_index);                      ④
    }
}

...
\end{codeboxt}
\noindent{\fontsize{9bp}{12bp}\selectfont ① 基例 1：立即返回。}\par
\noindent{\fontsize{9bp}{12bp}\selectfont ② 基例 2：必要时交换值。}\par
\noindent{\fontsize{9bp}{12bp}\selectfont ③ 获取基准元素的下标。}\par
\noindent{\fontsize{9bp}{12bp}\selectfont ④ 递归地对基准元素两侧的两个子数组排序。}\par

成员函数 \texttt{partition(\allowbreak{})\allowbreak{}} 实现划分一个子数组的各个步骤，该子数组最左边和最右边元素的下标分别是 \texttt{left\_\allowbreak{}index} 和 \texttt{right\_\allowbreak{}index}。

\begin{codeboxt}{代码清单 15.12（程序 15.5 Quicksort）：\texttt{Quicksort.\allowbreak{}cpp}（第 2 部分，共 3 部分）}
...

int Quicksort::partition(const int left_index, const int right_index)
{
    int middle_index = (left_index + right_index)/2;             ①
    int pivot_value  = data[middle_index];                       ①

    swap_values_at(middle_index, right_index);                   ②

    int i = left_index - 1;
    int j = right_index;

    while (i < j)
    {
        do
        {
            i++;                                                 ③
        } while ((i < right_index) && (data[i] < pivot_value));

        do
        {
            j--;                                                 ④
        } while ((j >= left_index) && (data[j] > pivot_value));

        if (i < j) swap_values_at(i, j);                         ⑤
    }

    swap_values_at(i, right_index);                              ⑥
    return i;
}
...
\end{codeboxt}
\noindent{\fontsize{9bp}{12bp}\selectfont ① 选取子数组中部附近的一个基准元素。}\par
\noindent{\fontsize{9bp}{12bp}\selectfont ② 把基准元素挪到子数组的右端。}\par
\noindent{\fontsize{9bp}{12bp}\selectfont ③ 把 i 向右移动。}\par
\noindent{\fontsize{9bp}{12bp}\selectfont ④ 把 j 向左移动。}\par
\noindent{\fontsize{9bp}{12bp}\selectfont ⑤ 交换第 i 个和第 j 个元素的值。}\par
\noindent{\fontsize{9bp}{12bp}\selectfont ⑥ 把挪开的基准元素移到它的正确位置。}\par

其余的成员函数负责交换值，并验证数组已正确排序。重载的输出运算符 \texttt{<<} 打印该数组。

\begin{codeboxt}{代码清单 15.13（程序 15.5 Quicksort）：\texttt{Quicksort.\allowbreak{}cpp}（第 3 部分，共 3 部分）}
...

void Quicksort::swap_values_at(const int index1, const int index2)
{
    int temp = data[index1];
    data[index1] = data[index2];
    data[index2] = temp;
}

bool Quicksort::verify_sorted()
{
    for (int i = 0; i < size - 1; i++)
    {
        if (data[i] > data[i+1]) return false;
    }

    return true;
}

ostream& operator <<(ostream& ostr, const Quicksort& q)
{
    for (int i = 0; i < q.size; i++) ostr << setw(3) << q.data[i];

    return ostr;
}
\end{codeboxt}

测试程序生成一个随机取值数组用于快速排序。

\begin{codeboxt}{代码清单 15.14（程序 15.5 Quicksort）：\texttt{tester.\allowbreak{}cpp}}
#include <iostream>
#include <stdlib.h>
#include <time.h>

#include "Quicksort.h"

using namespace std;

int main()
{
    const int SIZE = 20;
    int data[SIZE];

    srand(time(0));
    for (int i = 0; i < SIZE; i++) data[i] = rand()%100;

    Quicksort qsorter(data, SIZE);
    cout << "Before:" << qsorter << endl;
    qsorter.sort();

    cout << " After:" << qsorter << endl;

    cout << (qsorter.verify_sorted() ? "SORTED" : "ERROR") << endl;
    return 0;
}
\end{codeboxt}

示例输出为

\begin{codebox}
Before: 69 38 40  5 17 69 60 26 53 36 24 10 67 64 80 83 18 10 68 97
 After:  5 10 10 17 18 24 26 36 38 40 53 60 64 67 68 69 69 80 83 97
SORTED
\end{codebox}

快速排序程序对 \emph{N} 个值排序时会进行多少次递归调用？每次划分操作都把一个基准值放到其正确位置。完成基准划分后，可能产生两次递归调用。因此，递归调用最多可达 2\emph{N} 次。幸运的是，那些无需递归调用就能把值放到正确位置的基例，会显著减少调用的总数。

\section{斐波那契数列与一场递归灾难}

递归是一种非常强大的工具，但和所有强力工具一样，我们必须谨慎使用。用递归进行不当设计时可能出问题的经典例子就是斐波那契数列（Fibonacci sequence）0、1、2、3、5、8、13、21……其中，从 0 和 1 开始，后续每一项都是前两项之和。该数列的数学定义为e.

\myGraphic{0.267}{images/CH15_F12_Mak_f1.png}{}

我们可能会忍不住照搬这个递归的数学定义，写出一个递归函数 \texttt{f()} 。

\begin{codeboxt}{代码清单 15.15（程序 15.6 Fibonacci）：\texttt{fibrecursive.\allowbreak{}cpp}（递归灾难）}
#include <iostream>
#include <iomanip>

long f(const int n);

using namespace std;

int main()
{
    cout << «  n        f[n]» << endl;
    cout << «---------------» << endl;

    for (int n = 0; n < 50; n++)
    {
        cout << setw(3) << n << setw(12) << f(n) << endl;
    }

    return 0;
}

long f(const int n)
{
    if      (n == 0) return 0;
    else if (n == 1) return 1;
    else             return f(n-2) + f(n-1);
}
\end{codeboxt}

输出一开始

\begin{codebox}
  n        f[n]
---------------
  0           0
  1           1
  2           1
  3           2
  4           3
  5           5
  6           8
  7          13
  8          21
  9          34
 10          55
 11          89
 12         144
 13         233
 14         377
 15         610
 16         987
 17        1597
 18        2584
 19        4181
 20        6765
 21       10946
 22       17711
 23       28657
 24       46368
 25       75025
 26      121393
 27      196418
 28      317811
 29      514229
 30      832040
 31     1346269
 32     2178309
 33     3524578
 34     5702887
 35     9227465
 36    14930352
 37    24157817
 38    39088169
 39    63245986
 40   102334155
...
\end{codebox}

但会变得越来越慢。图 15.13 说明了原因。

\myGraphic{0.831}{images/CH15_F13_Mak.png}{图 15.13 计算 \texttt{f(6)} 的递归调用树。其中已经出现了许多参数相同的重复调用。随着 \texttt{n} 变大，计算 \texttt{f(n)} 的这棵树会越长越大，计算耗时将呈指数级增长。}

递归解可能看起来简单而优雅，但如果我们不理解递归调用是如何发生的、或者一共有多少次调用，就可能遭遇非常糟糕的性能意外。所有函数调用，无论是否递归，都会带来运行时和内存开销。

\section{动态回溯增强递归的能力}

递归与动态回溯的结合是一种极其强大的设计工具，我们可以用它来求解涉及多个步骤、且每一步都有若干可能求解路径的编程问题。从起始步骤开始，程序用递归依次探索每一条求解路径——每次递归调用都让代码沿着某条路径处理一个规模更小的问题。当问题小到能立即看出是否找到了解时，就是基例，此时递归调用分别返回 true 或 false。

在每一步，如果某次递归调用返回 false，意味着它正在探索的求解路径走进了死胡同、以失败告终，代码就回溯，并递归地探索下一条求解路径。在探索完某一步的所有路径之后，程序进一步回溯到上一步，在那里继续递归地探索其他的求解路径。当程序一路回溯到起始步骤、并探索完该步骤的所有求解路径时，程序结束。带动态回溯的递归往往构成问题的暴力解法。我们利用计算机的速度来尝试众多求解路径。

我们的第一个示例程序用带回溯的递归求解八皇后问题。第二个示例程序求解数独谜题。

\section{用递归与回溯求解八皇后问题}

这个谜题的目标是在棋盘上摆放八个皇后，使得任何两个皇后都不能互相攻击，无论是水平、垂直还是对角线方向。一个示例解是

\begin{codebox}
Q . . . . . . . 
. . . . . . Q . 
. . . . Q . . . 
. . . . . . . Q 
. Q . . . . . . 
. . . Q . . . . 
. . . . . Q . . 
. . Q . . . . . 
\end{codebox}

把镜像和旋转的情况都算上，共有 92 个解。

类 \texttt{Queens} 每次处理一列来解这个谜题，从最左边的一列开始（代码清单 15.16），每一列构成解的一个步骤。在每一列中，程序为一个皇后寻找安全位置。某一列中皇后的安全位置，是指该皇后不会被已经放置在左侧各列中的任何其他皇后攻击到的位置。如果找到安全位置，程序就在那里放置一个皇后，然后调用成员函数 \texttt{find\_\allowbreak{}solutions(\allowbreak{})\allowbreak{}}，通过递归地测试下一列来尝试解出谜题。这个子问题更小，因为需要考察的剩余列少了一列，我们离最后一列这个基例也更近了。

从某一列中皇后的位置出发进行的递归求解搜索，最终要么成功，要么失败。如果找到了一个解，程序就打印它。无论哪种情况，程序都在该列中把皇后下移到下一个安全位置。移动皇后就是动态回溯，它为从新位置递归搜索另一个解做好了准备。如果某一列中再没有安全位置，程序就回溯到上一列——即退回一步——并把那里的皇后移到该列的下一个安全位置。

基例是尝试把皇后放到最后一列中的某个安全位置，这可能成功也可能失败。当程序回溯到第一列、并尝试过该列的每个位置时，程序结束。私有布尔矩阵 \texttt{occupied} 记录各皇后的位置。

\begin{codeboxt}{代码清单 15.16（程序 15.7 Queens）：\texttt{Queens.h}}
const int SIZE = 8;

class Queens
{
public:
    Queens() : count(0)
    {
        for (int row = 0; row < SIZE; row ++)
        {
            for (int col = 0; col < SIZE; col ++)
            {
                occupied[row][col] = false;
            }
        }
    }

    void print_solutions();

private:
    bool occupied[SIZE][SIZE];                               ①
    int count;

    void find_solutions(const int col);                      ②
    bool is_safe_position(const int row, const int col) const;
    void print();
};
\end{codeboxt}
\noindent{\fontsize{9bp}{12bp}\selectfont ① 各皇后的位置}\par
\noindent{\fontsize{9bp}{12bp}\selectfont ② 递归函数}\par

公有成员函数 \texttt{print\_\allowbreak{}solutions(\allowbreak{})\allowbreak{}} 针对起始列 0 调用私有递归成员函数 \texttt{find\_\allowbreak{}solutions(\allowbreak{})\allowbreak{}}（代码清单 15.17）。该函数逐个检查某一列的每个位置是否为安全位置。如果在该列中找到了一个皇后的安全位置，它就递归地调用自身，通过测试下一列来推进到求解路径的下一步。

在递归地检查完某一列中的每个安全位置之后，函数通过把皇后移到该列的下一个安全位置来回溯，以便从那里继续尝试新的求解路径。如果该列再没有安全位置，\texttt{find\_\allowbreak{}solutions(\allowbreak{})\allowbreak{}} 就返回，求解搜索回溯到上一列，从那里继续尝试求解路径。最后一列是基例——不会从该列发出任何递归调用。

\begin{codeboxt}{代码清单 15.17（程序 15.7 Queens）：\texttt{Queens.\allowbreak{}cpp}（第 1 部分，共 2 部分）}
#include <iostream>
#include "Queens.h"

using namespace std;

void Queens::print_solutions()
{
    find_solutions(0);                                    ①
}

void Queens::find_solutions(const int col)
{
    for (int row = 0; row < SIZE; row++)                  ②
    {
        if (is_safe_position(row, col))
        {
            occupied[row][col] = true;                    ③

            if (col == SIZE - 1) print();                 ④
            else                 find_solutions(col + 1); ⑤
        }

        occupied[row][col] = false;                       ⑥
    }
}
\end{codeboxt}
\noindent{\fontsize{9bp}{12bp}\selectfont ① 针对列 0 的初始调用}\par
\noindent{\fontsize{9bp}{12bp}\selectfont ② 循环回溯到本列的下一个安全位置。}\par
\noindent{\fontsize{9bp}{12bp}\selectfont ③ 在一个安全位置放置皇后。}\par
\noindent{\fontsize{9bp}{12bp}\selectfont ④ 基例：在最后一列找到了一个解。}\par
\noindent{\fontsize{9bp}{12bp}\selectfont ⑤ 递归调用，以在下一列中寻找一个解}\par
\noindent{\fontsize{9bp}{12bp}\selectfont ⑥ 回溯：移除皇后，尝试本列的下一个安全位置。}\par

\myGraphic{0.849}{images/CH15_F13_UN09_Mak.png}{}

私有布尔成员函数 \texttt{is\_\allowbreak{}safe\_\allowbreak{}position(\allowbreak{})\allowbreak{}} 检查棋盘上的某个位置是否会受到前面几列中已放置皇后的攻击，如下面的代码清单所示。它检查皇后左侧的同一行、从皇后到左上方的对角线，以及从皇后到左下方的对角线。由于每一列只有一个皇后，因此不可能受到同一列中另一个皇后的攻击。私有成员函数 \texttt{print()} 打印一个解。

\begin{codeboxt}{代码清单 15.18（程序 15.7 Queens）：\texttt{Queens.\allowbreak{}cpp}（第 2 部分，共 2 部分）}
bool Queens::is_safe_position(const int row, const int col) const
{
    int c = col - 1;

    while (c >= 0)                           ①
    {                                        ①
        if (occupied[row][c]) return false;  ①
        c--;                                 ①
    }                                        ①

    int r = row - 1;
    c = col - 1;

    while ((r >= 0) && (c >= 0))             ②
    {                                        ②
        if (occupied[r][c]) return false;    ②
        r--;                                 ②
        c--;                                 ②
    }                                        ②

    r = row + 1;
    c = col - 1;

    while ((r < SIZE) && (c >= 0))           ③
    {    #C                                  ③
        if (occupied[r][c]) return false;    ③
        r++;                                 ③
        c--;                                 ③
    }                                        ③

    return true;
}

void Queens::print()
{
    count++;

    cout << endl << "Solution #" << count << endl << endl;

    for (int row = 0; row < SIZE; row++)
    {
        for (int col = 0; col < SIZE; col++)
        {
            cout << (occupied[row][col] ? "Q " : ". ");
        }
        cout << endl;
    }
}
\end{codeboxt}
\noindent{\fontsize{9bp}{12bp}\selectfont ① 检查皇后左侧的同一行。}\par
\noindent{\fontsize{9bp}{12bp}\selectfont ② 检查从皇后到左上方的对角线。}\par
\noindent{\fontsize{9bp}{12bp}\selectfont ③ 检查从皇后到左下方的对角线。}\par

测试程序打印这些解。

\begin{codeboxt}{代码清单 15.19（程序 15.7 Queens）：\texttt{tester.\allowbreak{}cpp}}
#include <iostream>
#include "Queens.h"

using namespace std;

int main()
{
    Queens q;
    q.print_solutions();

    cout << endl << "Done!" << endl;
    return 0;
}
\end{codeboxt}

输出包含全部 92 个解：

\begin{codebox}
Solution #1

Q . . . . . . . 
. . . . . . Q . 
. . . . Q . . . 
. . . . . . . Q 
. Q . . . . . . 
. . . Q . . . . 
. . . . . Q . . 
. . Q . . . . . 

Solution #2

Q . . . . . . . 
. . . . . . Q . 
. . . Q . . . . 
. . . . . Q . . 
. . . . . . . Q 
. Q . . . . . . 
. . . . Q . . . 
. . Q . . . . . 

Solution #3

Q . . . . . . . 
. . . . . Q . . 
. . . . . . . Q 
. . Q . . . . . 
. . . . . . Q . 
. . . Q . . . . 
. Q . . . . . . 
. . . . Q . . . 

...

Solution #92

. . Q . . . . . 
. . . . . Q . . 
. . . Q . . . . 
. Q . . . . . . 
. . . . . . . Q 
. . . . Q . . . 
. . . . . . Q . 
Q . . . . . . . 
Done!
\end{codebox}

\section{用递归与回溯求解数独谜题}

我们可以用同样“递归加回溯”的策略设计一个求解数独谜题的程序。例如，下面是一个输入文件，其中包含给定的格子值，其余位置为 0：

\begin{codebox}
0 0 0   0 0 0   0 2 0
3 0 0   0 0 6   0 0 0
0 0 0   0 8 0   4 7 0

2 0 0   3 0 0   0 0 0
0 1 0   0 4 0   0 8 0
0 0 0   0 0 5   0 0 6

0 8 7   0 1 0   0 0 0
0 0 0   9 0 0   0 0 5
0 4 0   0 0 0   0 0 0
\end{codebox}

该程序产生如下输出：

\begin{codebox}
Input:

. . . | . . . | . 2 . 
3 . . | . . 6 | . . . 
. . . | . 8 . | 4 7 . 
------+-------+-------
2 . . | 3 . . | . . . 
. 1 . | . 4 . | . 8 . 
. . . | . . 5 | . . 6 
------+-------+-------
. 8 7 | . 1 . | . . . 
. . . | 9 . . | . . 5 
. 4 . | . . . | . . . 

Solution:

8 5 9 | 7 3 4 | 6 2 1 
3 7 4 | 1 2 6 | 5 9 8 
6 2 1 | 5 8 9 | 4 7 3 
------+-------+-------
2 6 8 | 3 9 1 | 7 5 4 
7 1 5 | 6 4 2 | 3 8 9 
4 9 3 | 8 7 5 | 2 1 6 
------+-------+-------
5 8 7 | 4 1 3 | 9 6 2 
1 3 2 | 9 6 7 | 8 4 5 
9 4 6 | 2 5 8 | 1 3 7 
\end{codebox}

为求解一个数独谜题，程序从左上角开始，逐格移动，从左到右、自上而下，直到右下角。每个格子是解的一个步骤。谜题开始时，部分格子里已经填好了给定的数字。在每个未填的格子处，程序首先尝试数字 1，然后递归地尝试从下一个未填格子开始填完其余格子。这些递归尝试要么全部成功地填完剩余格子、从而找到了解，要么在途中某处失败。对下一个格子进行递归尝试是一个更小的问题，因为剩余的格子少了一个。基例是右下角的最后一个格子。

如果某个格子上的数字失败，程序就回溯，为那个格子尝试下一个数字。如果某个格子上九个数字都失败了，程序就进一步回溯，退回到上一个格子，并为那个格子尝试下一个数字。

类 \texttt{Sudoku} 有一个私有递归成员函数 \texttt{solution\_\allowbreak{}found(\allowbreak{})\allowbreak{}}，它返回 true 或 false，表示在某个格子中填入一个数字后是否找到了解。私有布尔函数 \texttt{is\_\allowbreak{}number\_\allowbreak{}ok(\allowbreak{})\allowbreak{}} 测试某个数字是否允许放入某个格子。

\begin{codeboxt}{代码清单 15.20（程序 15.8 Sudoku）：\texttt{Sudoku.h}}
class Sudoku
{
public:
    Sudoku(const char *file_name);

    bool solve() { return has_solution(0, 0); }

    void print() const;

private:
    static const int BLOCK_SIZE = 3;
    static const int GRID_SIZE  = BLOCK_SIZE*BLOCK_SIZE;

    int grid[GRID_SIZE][GRID_SIZE];                              ①

    bool has_solution(int row, int col);                         ②
    bool is_number_ok(int row, int col, int number) const;
};
\end{codeboxt}
\noindent{\fontsize{9bp}{12bp}\selectfont ① 谜题网格}\par
\noindent{\fontsize{9bp}{12bp}\selectfont ② 递归函数}\par

构造函数读取输入文件，用给定的值初始化网格，其余位置为 0。

\begin{codeboxt}{代码清单 15.21（程序 15.8 Sudoku）：\texttt{Sudoku.\allowbreak{}cpp}（第 1 部分，共 3 部分）}
#include <iostream>
#include <fstream>
#include "Sudoku.h"

using namespace std;

Sudoku::Sudoku(const char *file_name)
{
    try
    {
        ifstream grid_file(file_name);
        grid_file.exceptions(ifstream::failbit | ifstream::badbit);

        for (int row = 0; row < GRID_SIZE; row++)
        {
            for (int col = 0; col < GRID_SIZE; col++)
            {
                grid_file >> grid[row][col];
            }
        }
    }
    catch (ifstream::failure& ex)
    {
        cerr << "Input file error: " << file_name << endl;
        exit(-1);
    }
}

...
\end{codeboxt}

在把某个数字填入由给定 \texttt{row} 和 \texttt{col} 指定的网格格子之后，布尔成员函数 \texttt{solution\_\allowbreak{}found(\allowbreak{})\allowbreak{}} 递归地测试后续格子，检查是否找到了解。\texttt{for} 循环控制回溯。如果它填入的数字无法得出解，循环就尝试下一个数字。

\begin{codeboxt}{代码清单 15.22（程序 15.8 Sudoku）：\texttt{Sudoku.\allowbreak{}cpp}（第 2 部分，共 3 部分）}
...

bool Sudoku::solution_found(int row, int col)
{
    if ((row == GRID_SIZE-1) && (col == GRID_SIZE)) return true;  ①
    if (col == GRID_SIZE)                                         ②
    {
        row++;
        col = 0;
    }

    if (grid[row][col] != 0) return solution_found(row, col + 1); ③

    for (int number = 1; number <= GRID_SIZE; number++)           ④
    {
        if (is_number_ok(row, col, number))
        {
            grid[row][col] = number;                              ⑤
            if (solution_found(row, col + 1)) return true;        ⑥
        }
    }

    grid[row][col] = 0;                                           ⑦
    return false;                                                 ⑧
}
...
\end{codeboxt}
\noindent{\fontsize{9bp}{12bp}\selectfont ① 基例：到达最后一个格子并找到了解}\par
\noindent{\fontsize{9bp}{12bp}\selectfont ② 回绕到下一行的第一个格子。}\par
\noindent{\fontsize{9bp}{12bp}\selectfont ③ 如果该格子包含的是给定的数字，则递归地测试下一个格子。}\par
\noindent{\fontsize{9bp}{12bp}\selectfont ④ 循环测试从 1 开始的每个数字。}\par
\noindent{\fontsize{9bp}{12bp}\selectfont ⑤ 试探性地把一个允许的数字填入该格子。}\par
\noindent{\fontsize{9bp}{12bp}\selectfont ⑥ 如果该数字导致找到了解，就返回 true。否则回溯并尝试下一个数字。}\par
\noindent{\fontsize{9bp}{12bp}\selectfont ⑦ 在尝试完所有数字仍未成功之后，把该格子重置为“空”。}\par
\noindent{\fontsize{9bp}{12bp}\selectfont ⑧ 回溯到上一个格子：返回 false，以便上一次调用可以尝试下一个数字。}\par

这里有两个递归调用。第一个在当前格子包含的是给定数字时，递归地尝试下一个格子。在 \texttt{for} 循环中，把一个允许的数字赋给某个格子之后，第二个递归调用检查下一个格子，看从那里是否能找到解。如果找到了解，函数调用就返回 true。否则，循环尝试下一个数字。在尝试完所有数字仍未成功之后，函数把该格子重置为 0（“空”）并返回 false。这会使 \texttt{solution\_\allowbreak{}found(\allowbreak{})\allowbreak{}} 的调用者回溯，为上一个格子尝试下一个允许的数字。

终止递归的基例是在右下角的最后一个格子中放入一个允许的数字。回溯到上一个格子发生在某个格子中所有数字都尝试失败之后。当所有递归调用都展开回到第一个格子时，程序要么找到了一个解，要么这个谜题无解。

布尔成员函数 \texttt{is\_\allowbreak{}number\_\allowbreak{}ok(\allowbreak{})\allowbreak{}} 检查某个数字是否允许放入某个格子。它检查该数字是否已经出现在该格子所在的行、列或宫中。成员函数 \texttt{print()} 打印解。

\begin{codeboxt}{代码清单 15.23（程序 15.8 Sudoku）：\texttt{Sudoku.\allowbreak{}cpp}（第 3 部分，共 3 部分）}
bool Sudoku::is_number_ok(int row, int col, int number) const
{
    for (int c = 0; c < GRID_SIZE; c++)                       ①
    {                                                         ①
        if (grid[row][c] == number) return false;             ①
    }                                                         ①

    for (int r = 0; r < GRID_SIZE; r++)                       ②
    {                                                         ②
        if (grid[r][col] == number) return false;             ③
    }                                                         ③

    int block_row_start = row - row%BLOCK_SIZE;
    int block_col_start = col - col%BLOCK_SIZE;
    int block_row_end = block_row_start + BLOCK_SIZE;
    int block_col_end = block_col_start + BLOCK_SIZE;

    for (int r = block_row_start; r < block_row_end; r++)     ④
    {                                                         ④
        for (int c = block_col_start; c < block_col_end; c++) ④
        {                                                     ④
            if (grid[r][c] == number) return false;           ④
        }                                                     ④
    }                                                         ④

    return true;
}

void Sudoku::print() const
{
    for (int row = 0; row < GRID_SIZE; row++)
    {
        if ((row > 0) && (row%BLOCK_SIZE == 0))
        {
            cout << "------+-------+-------" << endl;
        }

        for (int col = 0; col < GRID_SIZE; col++)
        {
            if ((col > 0) && (col%BLOCK_SIZE == 0)) cout << "| ";

            int g = grid[row][col];
            if (g > 0) cout << g << " ";
            else       cout << ". ";
        }

        cout << endl;
    }
}
\end{codeboxt}
\noindent{\fontsize{9bp}{12bp}\selectfont ① 检查该格子所在的行。}\par
\noindent{\fontsize{9bp}{12bp}\selectfont ② 检查该格子所在的列。}\par
\noindent{\fontsize{9bp}{12bp}\selectfont ③ 检查该格子所在的列。}\par
\noindent{\fontsize{9bp}{12bp}\selectfont ④ 检查该格子所在的宫。}\par

测试程序打印带有给定数字的初始网格，调用类 \texttt{Sudoku} 的 \texttt{solve()} 成员函数，并在有解时打印解网格。

\begin{codeboxt}{代码清单 15.24（程序 15.8 Sudoku）：\texttt{tester.\allowbreak{}cpp}}
#include <iostream>
#include "Sudoku.h"

using namespace std;

int main(int argc, char *argv[])
{
    Sudoku sudoku(argv[1]);
    cout << "Input:" << endl << endl;
    sudoku.print();

    bool solved = sudoku.solve();                 ①

    if (solved)
    {
        cout << endl << «Solution:» << endl << endl;
        sudoku.print();
    }
    else
    {
        cout << endl << "No solution" << endl;
    }

    return 0;
}
\end{codeboxt}
\noindent{\fontsize{9bp}{12bp}\selectfont ① 对递归函数 solve() 的初始调用}\par

\myGraphic{0.746}{images/CH15_F13_UN10_Mak.png}{}

\section*{小结}

\begin{itemize}
\item 递归是一种软件设计技术，它通过把一个问题分解为一个更小但相似的子问题来求解。我们用同样的方法递归地求解每个子问题。
\item 递归在基例处停止，此时问题小到有一个显而易见、唾手可得的解。然后，随着递归展开，原问题得以求解。
\item 如果设计得当，递归的编程解可以既简单又优雅。然而，使用不当会导致严重的性能问题。
\item 对于需要探索多条求解路径的程序，递归与动态回溯的结合是一种强有力的设计工具。
\item 每一步都可能有若干条求解路径，用递归来探索它们。尝试完一条路径后，就回溯并尝试下一条路径。
\item 当某一步的所有求解路径都已尝试完毕时，回溯到上一步，并继续尝试上一步的各条路径。
\end{itemize}
